\abstract

\keywords{ПОИСК ПУТИ В ГРАФЕ, МНОГОАГЕНТНОЕ ПЛАНИРОВАНИЕ, АВТОНОМНЫЕ АГЕНТЫ, 
АДАПТИВНЫЙ ПОИСК ПУТИ, РАЗРЕШЕНИЕ КОНФЛИКТОВ}

Объектом исследования в данной работе является централизованное
управление группой автономных агентов, движущихся по графу
дорожной сети с узкими местами.

Цель работы --- разработать методы локального разрешения
конфликтов для глобального управления автономными агентами 
на графе дорожной сети с узкими местами и сопоставить их с классическими
методами многоагентного планирования.

Для достижения поставленной цели были проведены исследования в области
алгоритмов многоагентного планирования (CBS, CCBS, LaCAM, WHCA*), методов
одноагентного поиска путей (A*, SIPP) и методов разрешения конфликтов на
узких участках дорожной сети.

Основное содержание работы состояло в разработке двух методов локального
разрешения конфликтов (реактивный разъезд и система семафоров), адаптивной
модификации A* с эмпирической оценкой скоростей рёбер, а также модульного
симулятора городской среды для воспроизводимого сравнения алгоритмов.

Основными результатами работы являются предложенные алгоритмы управления,
программная среда симуляции и экспериментальное сравнение с оптимальным
алгоритмом CCBS на фрагментах реальной городской карты, показавшее, что
сочетание адаптивного A* и локального разрешения конфликтов не уступает
многоагентному планированию, а на типичных размерах задачи превосходит его.

Результаты разработки предназначены для использования при проектировании
систем централизованного управления парками автономных доставочных агентов.

Использование результатов данной работы позволяет существенно снизить
вычислительные затраты на согласование движения большого числа агентов при
сохранении сопоставимого качества совместного движения.
